\section{Grundlegende Rechenoperationen}
\label{sec:vector_algebra_operations}

Ähnlich wie für Zahlen sind auch für Vektoren verschiedene Rechenoperationen definiert. Sie
unterscheiden sich darin, ob sie Vektoren oder Skalare verknüpfen und ob ihr Ergebnis ein
Vektor oder ein Skalar ist. Wir beginnen mit einem Überblick über alle Rechenoperationen:

\begin{center}
  \renewcommand{\arraystretch}{1.35}
  \begin{tabularx}{0.96\linewidth}{@{}
      >{\raggedright\arraybackslash}X
      >{\centering\arraybackslash}p{3.2cm}
      >{\centering\arraybackslash}p{5.6cm}@{}}
    \toprule
    \textbf{Operation} & \textbf{Notation} & \textbf{Abbildung} \\
    \midrule
    Addition
      & $\vec a+\vec b$
      & Vektor $\times$ Vektor $\to$ Vektor \\
    Subtraktion
      & $\vec a-\vec b$
      & Vektor $\times$ Vektor $\to$ Vektor \\
    Multiplikation mit einem Skalar
      & $\lambda\vec a$
      & Skalar $\times$ Vektor $\to$ Vektor \\
    Skalarprodukt
      & $\vec a\cdot\vec b$
      & Vektor $\times$ Vektor $\to$ Skalar \\
    Kreuzprodukt
      & $\vec a\times\vec b$
      & Vektor $\times$ Vektor $\to$ Vektor \\
    \bottomrule
  \end{tabularx}
\end{center}

In diesem Abschnitt werden zunächst Addition, Subtraktion und Multiplikation
mit einem Skalar rechnerisch eingeführt und geometrisch interpretiert. Das
Skalarprodukt folgt in Abschnitt~\ref{sec:vector_algebra_scalar_product}; das
Kreuzprodukt in Abschnitt~\ref{sec:vector_algebra_cross_product}.

\subsection{Addition von Vektoren}
\label{subsec:vector_algebra_addition}

\begin{DefinitionBox}{Vektoraddition}{vector_algebra_addition}
  Die \emph{Vektoraddition} \EnglishTerm{vector addition} zweier Vektoren
  $\vec a,\vec b\in\R^n$ erfolgt komponentenweise:
  \[
    \vec a+\vec b
    =\begin{pmatrix}a_1\\a_2\\\vdots\\a_n\end{pmatrix}
     +\begin{pmatrix}b_1\\b_2\\\vdots\\b_n\end{pmatrix}
    =\begin{pmatrix}
      \Cloze{2.4cm}{a_1+b_1}\\[0.25em]
      \Cloze{2.4cm}{a_2+b_2}\\[0.25em]
      \vdots\\[0.25em]
      \Cloze{2.4cm}{a_n+b_n}
    \end{pmatrix}.
  \]
  Vektoren unterschiedlicher Dimension können nicht addiert werden.
\end{DefinitionBox}

\begin{KeyPointBox}{Geometrische Vektoraddition}{vector_algebra_geometric_addition}
  Um die Summe zweier Vektoren geometrisch zu konstruieren, zeichnet man
  $\vec a$ und $\vec b$ mit gemeinsamem Anfangspunkt und ergänzt sie zu einem
  Parallelogramm. Die vom gemeinsamen Anfangspunkt ausgehende Diagonale ist der
  Summenvektor $\vec a+\vec b$.
  \begin{center}
    \CachedTikzFigure{content/resources/chapters/vector_algebra/tikz/vector_addition.tex}
  \end{center}
\end{KeyPointBox}

\begin{ExampleBox}{Vektoren rechnerisch und geometrisch addieren}{vector_algebra_addition_example_2d}
  \begin{minipage}[c]{0.52\linewidth}
    Gegeben seien die beiden Vektoren
    \[
      \vec a=\begin{pmatrix}1\\2\end{pmatrix},
      \qquad
      \vec b=\begin{pmatrix}3\\1\end{pmatrix}.
    \]
    Die komponentenweise Addition liefert
    \[
      \vec a+\vec b=
      \begin{ClozeContinuation}[12]{2}
        \begin{pmatrix}1+3\\2+1\end{pmatrix}
        =\begin{pmatrix}4\\3\end{pmatrix}
      \end{ClozeContinuation}
    \]
  \end{minipage}\hfill
  \begin{minipage}[c]{0.44\linewidth}
    \centering
    \CachedTikzFigure{content/resources/chapters/vector_algebra/tikz/vector_addition_example_2d.tex}
  \end{minipage}
  \par\smallskip
  Rechnerisch und geometrisch erhalten wir somit dasselbe Ergebnis.
\end{ExampleBox}

Da Vektoren komponentenweise addiert werden, lassen sich die Rechengesetze
der Addition von Zahlen unmittelbar auf die Vektoraddition übertragen.

\begin{TheoremBox}{Rechenregeln der Vektoraddition}{vector_algebra_addition_laws}
  Bei der Addition beliebiger Vektoren $\vec a,\vec b,\vec c\in\R^n$ gelten die
  üblichen algebraischen Rechenregeln:
  \[
    \begin{aligned}
      \vec a+\vec b
        &=\Cloze{5cm}{\vec b+\vec a}
        &&\qquad\text{(Kommutativgesetz)},\\[0.5em]
      \vec a+(\vec b+\vec c)
        &=\Cloze{5cm}{(\vec a+\vec b)+\vec c}
        &&\qquad\text{(Assoziativgesetz)}
    \end{aligned}
  \]
\end{TheoremBox}

Daher können Vektoren beim Umstellen von Gleichungen hinsichtlich Addition
und Subtraktion genauso wie Zahlen behandelt werden.

\subsubsection{Subtraktion als Addition des Gegenvektors}
\label{subsubsec:vector_algebra_subtraction}

Für die Subtraktion führen wir keine neue Rechenart ein. Sie baut auf der
Vektoraddition und dem Gegenvektor auf.

\begin{DefinitionBox}{Vektorsubtraktion}{vector_algebra_subtraction}
  Die \emph{Vektorsubtraktion} \EnglishTerm{vector subtraction} ist als
  Addition von $\vec a$ mit dem Gegenvektor $-\vec b$ definiert:
  \[
    \vec a-\vec b
    :=\vec a+(-\vec b)
    =\begin{pmatrix}a_1\\a_2\\\vdots\\a_n\end{pmatrix}
     +\begin{pmatrix}-b_1\\-b_2\\\vdots\\-b_n\end{pmatrix}
    =\begin{pmatrix}
      \Cloze{2.4cm}{a_1-b_1}\\[0.25em]
      \Cloze{2.4cm}{a_2-b_2}\\[0.25em]
      \vdots\\[0.25em]
      \Cloze{2.4cm}{a_n-b_n}
    \end{pmatrix}.
  \]
\end{DefinitionBox}

\begin{KeyPointBox}{Geometrische Vektorsubtraktion}{vector_algebra_geometric_subtraction}
  Gemäß der Definition schreiben wir die Differenz als Summe mit dem
  Gegenvektor:
  \[
    \vec a-\vec b=\vec a+(-\vec b).
  \]
  Zur geometrischen Konstruktion zeichnen wir, wie bei einer Summe üblich,
  $\vec a$ und $-\vec b$ mit gemeinsamem Anfangspunkt zu einem Parallelogramm.
  Dessen Diagonale ist der Differenzvektor $\vec a-\vec b$.
  \begin{center}
    \CachedTikzFigure{content/resources/chapters/vector_algebra/tikz/vector_subtraction_as_addition.tex}
  \end{center}
  Zugleich ist $\vec a-\vec b$ der Verbindungsvektor von der Spitze von
  $\vec b$ zur Spitze von $\vec a$.
\end{KeyPointBox}

\begin{ExampleBox}{Vektoren subtrahieren}{vector_algebra_componentwise_subtraction}
  \begin{minipage}[c]{0.52\linewidth}
    Gegeben seien die beiden Vektoren
    \[
      \vec a=\begin{pmatrix}4\\3\end{pmatrix},
      \qquad
      \vec b=\begin{pmatrix}1\\2\end{pmatrix}.
    \]
    Die Subtraktion liefert
    \[
      \vec a-\vec b=
      \begin{ClozeContinuation}[12]{2}
        \vec a+(-\vec b)
        =\begin{pmatrix}4-1\\3-2\end{pmatrix}
        =\begin{pmatrix}3\\1\end{pmatrix}
      \end{ClozeContinuation}
    \]
    Rechnerisch und geometrisch erhalten wir somit dasselbe Ergebnis.
  \end{minipage}\hfill
  \begin{minipage}[c]{0.44\linewidth}
    \centering
    \CachedTikzFigure{content/resources/chapters/vector_algebra/tikz/vector_subtraction_example_2d.tex}
  \end{minipage}
  \par\smallskip
\end{ExampleBox}

\subsection{Multiplikation mit einem Skalar}
\label{subsec:vector_algebra_scalar_multiplication}

Die erste Mulitplikation, die wir in der Vektorrechnung kennenlernen,
ist die Mulitplikation eines Skalars mit einem Vektor.

\begin{DefinitionBox}{Multiplikation mit einem Skalar}{vector_algebra_scalar_multiplication}
  Die \emph{Multiplikation mit einem Skalar} \EnglishTerm{scalar multiplication}
  erfolgt für $\lambda\in\R$ und $\vec a\in\R^n$ komponentenweise:
  \[
    \lambda\vec a
    =\begin{ClozeContinuation}[15]{5}
      \lambda\begin{pmatrix}a_1\\a_2\\\vdots\\a_n\end{pmatrix}
      =\begin{pmatrix}\lambda a_1\\\lambda a_2\\\vdots\\\lambda a_n\end{pmatrix}
    \end{ClozeContinuation}.
  \]
  Das Ergebnis $\lambda\vec a\in\R^n$ ist wieder ein Vektor.
\end{DefinitionBox}

\begin{ExampleBox}{Einen Vektor mit einem Skalar multiplizieren}{vector_algebra_scale_example}
  \begin{minipage}[c]{0.49\linewidth}
    Gegeben sei der Vektor
    \[
      {\color{black}
        \vec v=\begin{pmatrix}2\\1\end{pmatrix}}.
    \]
    Die Multiplikation mit dem Skalar $\lambda$ für
    \begin{itemize}[leftmargin=1.4em,itemsep=0.8em]
      \item {\color{olive!80!black}
        $\lambda=2$ entspricht einer \Cloze{2.8cm}{Streckung}:}
        \[
          {\color{olive!80!black}2\vec v=}
          \begin{ClozeContinuation}[11]{2}
            {\color{olive!80!black}
              2\begin{pmatrix}2\\1\end{pmatrix}
              =\begin{pmatrix}4\\2\end{pmatrix}}
          \end{ClozeContinuation}.
        \]
      \item {\color{THMBlue}
        $\lambda=-1$ liefert den \Cloze{3.1cm}{Gegenvektor}:}
        \[
          {\color{THMBlue}(-1)\vec v=}
          \begin{ClozeContinuation}[12]{2}
            {\color{THMBlue}
              (-1)\begin{pmatrix}2\\1\end{pmatrix}
              =\begin{pmatrix}-2\\-1\end{pmatrix}}
          \end{ClozeContinuation}.
        \]
    \end{itemize}
  \end{minipage}\hfill
  \begin{minipage}[c]{0.48\linewidth}
    \centering
    \CachedTikzFigure{content/resources/chapters/vector_algebra/tikz/scalar_multiplication_example_2d.tex}
  \end{minipage}
\end{ExampleBox}

Die geometrische Wirkung der Multiplikation mit einem Skalar lässt sich
folgendermaßen zusammenfassen:

\begin{KeyPointBox}{Geometrische Wirkung der Multiplikation mit einem Skalar}{vector_algebra_scalar_effect}
  \begin{center}
    \CachedTikzFigure{content/resources/chapters/vector_algebra/tikz/scalar_multiplication_cases.tex}
  \end{center}
  \[
    \text{Für }
    \begin{gathered}
      \lambda>0\\[0.5em]
      \lambda<0
    \end{gathered}
    \text{ ist der resultierende Vektor }\lambda\vec a\text{ }
    \begin{gathered}
      \text{parallel}\\[0.5em]
      \text{antiparallel}
    \end{gathered}
    \text{ zu }\vec a,
    \qquad\text{d.\,h.}\qquad
    \begin{gathered}
      \lambda\vec a\upuparrows\vec a,\\[0.5em]
      \lambda\vec a\updownarrows\vec a.
    \end{gathered}
  \]
\end{KeyPointBox}

Da in jeder Komponente ausschließlich Zahlen multipliziert werden, übertragen
sich die entsprechenden Rechengesetze der Zahlenmultiplikation auf die
Multiplikation eines Vektors mit Skalaren.

\begin{TheoremBox}{Rechenregeln der Multiplikation mit einem Skalar}{vector_algebra_scalar_laws}
  Für $\alpha,\beta\in\R$ und $\vec a,\vec b\in\R^n$ gilt
  \[
    \begin{array}{@{}c@{\qquad}l@{}}
      \text{--}
        & \text{Kommutativgesetz}\\[1.2em]
      \alpha(\beta\vec a)=(\alpha\beta)\vec a
        & \text{Assoziativgesetz}\\[1.2em]
      \left.
        \begin{aligned}
          (\alpha+\beta)\vec a&=\alpha\vec a+\beta\vec a,\\[0.4em]
          \alpha(\vec a+\vec b)&=\alpha\vec a+\alpha\vec b
        \end{aligned}
        \hspace{0.8em}
      \right\}
        & \text{Distributivgesetze}
    \end{array}
  \]
\end{TheoremBox}

Daher kann bei der Multiplikation eines Vektors mit Skalaren so gerechnet
werden, wie man es von Zahlen gewohnt ist.

